\section{Verificación}\label{sec:verificacion}

\subsection{Vectores de prueba}

Los modelos de punto fijo generan los vectores con los que se verifica el RTL. Cada archivo de texto tiene una línea por muestra de entrada con ocho enteros: índice, marca de fin de trama, entrada I/Q y las salidas esperadas de la rama temporal y de la rama en frecuencia (Q2.7). Cada trama empieza con historial en cero. Un test falla si algún archivo versionado deja de coincidir con el generador.

\begin{table}[h]
  \centering
  \caption{Casos de prueba.}\label{tab:vectores}
  \begin{tabular}{lrrl}
    \toprule
    Archivo & Tramas & Muestras & Qué prueba \\
    \midrule
    \texttt{impulso} & 1 & 32 & Respuesta al impulso $1-j$ \\
    \texttt{ceros} & 1 & 32 & Salida nula \\
    \texttt{extremos\_tiempo} & 1 & 16 & Salida temporal en $\pm\sum|h_q|$ \\
    \texttt{extremos\_frecuencia} & 9 & 144 & Máximo de \texttt{work} en cada etapa \\
    \texttt{bordes} & 19 & 202 & Tramas de 1 a 17, 24 y 25 muestras \\
    \texttt{qpsk\_s1} & 1 & 2048 & QPSK 2×, semilla 1 (regresión y actividad) \\
    \bottomrule
  \end{tabular}
\end{table}

Los cinco primeros (426 muestras) corren en cada integración; \texttt{qpsk\_s1} se genera y se usa en la regresión. Cada RTL se compara palabra por palabra con su columna de salida, por índice de resultado aceptado y no por ciclo, y la marca de fin de trama de la salida debe coincidir con la del vector.

\pendiente{Andrés (A06)}{Banco de pruebas común: un mismo testbench para las variantes de ambas ramas, con pausas en la entrada y contrapresión en la salida; resultados de la regresión.}

\pendiente{Andrés}{Integración continua: qué corre \texttt{make smoke} (modelos, vectores, simulación con Icarus) y \texttt{make full} (síntesis con Yosys).}
